% =====================================================================
%  Bölüm 3 — Önermeler Mantığı II: Değilleme, Çelişki ve Denklik
% =====================================================================
\chapter{Önermeler Mantığı II: Değilleme, Çelişki ve Denklik}\label{bol:degilleme}

\begin{bolumhedefleri}
\begin{itemize}
  \item \lc{True} ve \lc{False} önermelerini ve Curry--Howard'daki
        karşılıklarını,
  \item Değillemenin \Lean{}'de neden \lc{¬P := P → False} olarak
        tanımlandığını,
  \item Çelişkiden her şeyin çıktığını (\emph{ex falso quodlibet}) ve
        \leantac{exfalso}, \lc{absurd}, \leantac{contradiction} araçlarını,
  \item Denkliği (\lc{↔}), \lc{Iff.mp}, \lc{Iff.mpr} ve \leantac{rw} ile
        yeniden yazmayı,
  \item De Morgan yasalarının hangi yarısının yapıcı olarak ispatlanabildiğini
        ve neden ötekinin ispatlanamadığını.
\end{itemize}
\end{bolumhedefleri}

% ---------------------------------------------------------------------
\section{İki uç: True ve False}\label{sec:true-false}

Önermeler arasında iki özel önerme vardır. \lc{True} her zaman doğru olan,
\lc{False} ise hiçbir zaman doğru olmayan önermedir. Mantık kitaplarında
bunlar genellikle $\top$ (\emph{tepe}) ve $\bot$ (\emph{dip}) ile gösterilir.

\begin{center}
\begin{tabular}{c@{\hspace{5em}}c}
\begin{prooftree}
  \infer0[$\top$G]{\top}
\end{prooftree}
&
\begin{prooftree}
  \hypo{\bot}
  \infer1[$\bot$E]{P}
\end{prooftree}
\\[10pt]
\small\textcolor{ikincil}{\bfseries giriş:} \leantac{trivial}, \lc{True.intro} &
\small\textcolor{vurgu}{\bfseries eleme (ex falso):} \lc{False.elim}, \leantac{exfalso}
\end{tabular}
\end{center}

\noindent İki kural birbirinin aynadaki görüntüsü gibidir:

\begin{itemize}
  \item $\top$'un \textbf{giriş} kuralının hiç öncülü yoktur: \lc{True}'yu
        her zaman, hiçbir şey bilmeden kanıtlayabiliriz. Buna karşılık
        \textbf{eleme} kuralı yoktur: \lc{True}'yu bilmek bize hiçbir bilgi
        vermez.
  \item $\bot$'un \textbf{giriş} kuralı yoktur: \lc{False}'u kanıtlamanın
        doğrudan bir yolu yoktur. Buna karşılık \textbf{eleme} kuralı
        çok güçlüdür: \lc{False}'un bir kanıtı elimize geçerse, \emph{her
        önermeyi} kanıtlayabiliriz.
\end{itemize}

\leandosya[True'nun kanıtı]{Bolum03}{dogru}

\subsection{Ex falso quodlibet}

Latince ``yanlıştan ne istersen'' anlamına gelen bu ilke, $\bot$E
kuralının adıdır.

\leandosya[Ex falso]{Bolum03}{ex_falso}

\noindent Daha ilginç bir ispat, \leantac{cases} taktiğini kullanır:

\leandosya[Ex falso, durum ayırarak]{Bolum03}{ex_falso_cases}

Neden işe yarıyor? \Cref{sec:veya}'te \leantac{cases}'in bir veya
kanıtını \emph{iki} duruma ayırdığını görmüştük, çünkü \lc{P ∨ Q}'nun
kanıtı iki biçimde kurulabilir (\lc{Or.inl} ve \lc{Or.inr}). \lc{False}'un
kanıtı ise \emph{hiçbir} biçimde kurulamaz. Dolayısıyla \lc{h : False}
üzerinde durum ayırınca \emph{sıfır} durum çıkar ve kapatılacak hedef
kalmaz.

\begin{figure}[htbp]
\centering
\begin{tikzpicture}[
    tip/.style={draw=#1, fill=#1!8, thick, rounded corners=4pt,
                minimum width=2.6cm, minimum height=1.0cm,
                font=\leanfont\small, align=center},
    kur/.style={draw=#1, fill=white, thick, rounded corners=3pt,
                font=\leanfont\scriptsize, inner sep=3pt},
    ok/.style={-{Stealth[length=2.3mm]}, thick, #1},
  ]
  % Or
  \node[tip=leanframe] (or) at (0,0) {P ∨ Q};
  \node[kur=ikincil] (o1) at (-1.0,-1.5) {Or.inl};
  \node[kur=ikincil] (o2) at (1.0,-1.5) {Or.inr};
  \draw[ok=ikincil] (or) -- (o1); \draw[ok=ikincil] (or) -- (o2);
  \node[font=\baslikfont\scriptsize, text=black!60] at (0,-2.3) {2 kurucu → 2 durum};
  % True
  \node[tip=bittiframe] (tr) at (5,0) {True};
  \node[kur=bittiframe] (t1) at (5,-1.5) {True.intro};
  \draw[ok=bittiframe] (tr) -- (t1);
  \node[font=\baslikfont\scriptsize, text=black!60] at (5,-2.3) {1 kurucu → 1 durum};
  % False
  \node[tip=hataframe] (fa) at (10,0) {False};
  \node[text=hataframe] at (10,-1.5) {\Large$\emptyset$};
  \node[font=\baslikfont\scriptsize, text=black!60] at (10,-2.3) {0 kurucu → 0 durum};
\end{tikzpicture}
\caption{Bir önermenin \emph{kurucuları}, kanıtını inşa etmenin yollarıdır.
\leantac{cases} her kurucu için bir durum açar. \lc{False}'un kurucusu
yoktur.}
\label{fig:kurucular}
\end{figure}

\begin{onemlifikir}[Önemli Fikir: True ve False birer tiptir]
Curry--Howard eşlemesinde \lc{True}, tek elemanlı tipe (programlamada
\emph{unit}), \lc{False} ise hiç elemanı olmayan \emph{boş tipe} karşılık
gelir. Boş tipten başka bir tipe giden fonksiyonu yazmak kolaydır: Hiçbir
girdi gelmeyeceği için hiçbir çıktı üretmemiz de gerekmez. Ex falso tam
olarak budur. Bu tipler \cref{bol:tumevarimsal-tipler}'da \leankw{inductive}
ile kendimiz tanımlayacağız.
\end{onemlifikir}

% ---------------------------------------------------------------------
\section{Değilleme: \texorpdfstring{$\neg P$}{¬P} ``\texorpdfstring{$P$}{P} ise çelişki'' demektir}\label{sec:degilleme}

\Lean{}'de değilleme yeni bir temel kavram değildir; implikasyon ve
\lc{False} ile \emph{tanımlanır}:

\leandosya[Değillemenin tanımı]{Bolum03}{not_tanim}

\begin{tanim}{Değilleme}{degilleme}
Bir $P$ önermesinin \textbf{değillemesi},
\[
  \neg P \;:=\; P \to \bot
\]
önermesidir. Yani ``$P$ değil'' demek, ``$P$ doğru olsaydı bir çelişkiye
varırdık'' demektir.
\end{tanim}

Bu tanımın önemli bir sonucu var: Değilleme için ayrı kurallar
öğrenmemize gerek yok, çünkü $\neg$ aslında $\to$'dur. $\neg$G kuralı
$\to$G'nin, $\neg$E kuralı da $\to$E'nin özel hâlidir:

\begin{center}
\begin{tabular}{c@{\hspace{4em}}c}
\begin{prooftree}
  \hypo{[P]^{1}}
  \ellipsis{}{\bot}
  \infer1[$\neg$G$^{1}$]{\neg P}
\end{prooftree}
&
\begin{prooftree}
  \hypo{\neg P}
  \hypo{P}
  \infer2[$\neg$E]{\bot}
\end{prooftree}
\\[10pt]
\small\textcolor{ikincil}{\bfseries giriş:} \leantac{intro hp}, hedef \lc{False} olur &
\small\textcolor{vurgu}{\bfseries eleme:} \lc{hnp hp : False}
\end{tabular}
\end{center}

\subsection{Çelişmezlik ilkesi}

Aristoteles'ten beri mantığın temel yasalarından biri sayılan
\emph{çelişmezlik ilkesi}, bir önermenin hem doğru hem yanlış
olamayacağını söyler: $\neg(P \wedge \neg P)$.

\leandosya[Çelişmezlik ilkesi]{Bolum03}{degil_giris}

\begin{hedefgecisi}
\begin{hedefdurumu}[\texttt{intro h} sonrası]
P : Prop
h : P ∧ ¬P
⊢ False
\end{hedefdurumu}
\sonra
\begin{hedefdurumu}[\texttt{obtain} sonrası]
P : Prop
hp : P
hnp : ¬P
⊢ False
\end{hedefdurumu}
\end{hedefgecisi}

\leantac{intro h} adımına dikkat edin: Hedef \lc{¬(P ∧ ¬P)} idi ve
görünürde bir ok yoktu. Ama \lc{¬X}, \lc{X → False}'un kısaltması
olduğu için \leantac{intro} onu açabildi ve hedef \lc{False} oldu.
Son adımda \lc{hnp : P → False} fonksiyonunu \lc{hp} argümanıyla
çağırarak \lc{False}'u elde ettik.

\leandosya[Çelişmezlik, terim modu]{Bolum03}{degil_giris_terim}

\subsection{Çift değillemenin girişi}

\leandosya[\texorpdfstring{$P \to \neg\neg P$}{P → ¬¬P}]{Bolum03}{cift_degil_giris}

\noindent Bu ispatı açık açık yazarsak:
\[
  P \to \neg\neg P
  \;=\; P \to (\neg P \to \bot)
  \;=\; P \to ((P \to \bot) \to \bot).
\]
\leantac{intro hp hnp} iki oku birden açar. Bu teoremin \emph{tersi},
yani $\neg\neg P \to P$, bu bölümde ispatlayamayacağımız ilk önemli
önermedir.

\begin{dikkat}[Dikkat: $\neg\neg P$ ile $P$ aynı şey değildir]
Okulda öğrendiğimiz mantıkta ``değilin değili'' kendisidir. \Lean{}'in
varsayılan mantığında ise $P$'den $\neg\neg P$ çıkar, ama tersi
çıkmaz. $\neg\neg P$ yalnızca ``$P$'nin yanlış olduğunu gösteremezsin''
der; $P$'nin bir kanıtını \emph{vermez}. Bu farkın ne kadar derin
olduğunu \cref{bol:yapici-klasik}'te göreceğiz.
\end{dikkat}

% ---------------------------------------------------------------------
\section{Çelişkiyle çalışmanın üç yolu}\label{sec:celiski}

Bağlamda hem \lc{hp : P} hem \lc{hnp : ¬P} varsa, hedef ne olursa olsun
ispat bitmiştir. \Lean{} bunu söylemek için üç farklı araç sunar:

\leandosya[Aynı ispat, üç araçla]{Bolum03}{celiski_uc_yol}

\begin{center}
\renewcommand{\arraystretch}{1.45}
\rowcolors{2}{tablosatir}{white}
\begin{tabularx}{\linewidth}{>{\raggedright\arraybackslash}p{3.0cm} X}
\rowcolor{tabloust}
\tablobaslik{Araç} & \tablobaslik{Ne zaman kullanılır?} \\
\leantac{exfalso}       & Hedefi \lc{False}'a çevirmek istediğinizde. Ardından çelişkiyi kendiniz kurarsınız. \\
\lc{absurd hp hnp}      & Çelişkili iki terimi doğrudan elinizde tuttuğunuzda (terim modu). \\
\leantac{contradiction} & Bağlamda açık bir çelişki olduğunda: \lc{P} ve \lc{¬P}, \lc{h : False}, ya da \lc{h : a ≠ a} gibi. \\
\end{tabularx}
\end{center}

\leantac{contradiction} en kısa olanıdır, ama yalnızca çelişki bağlamda
\emph{açıkça} duruyorsa çalışır. Çelişkiyi elde etmek için birkaç adım
gerekiyorsa \leantac{exfalso} ile başlamak daha açık bir ispat verir.

% ---------------------------------------------------------------------
\section{Kontrapozisyon ve modus tollens}\label{sec:kontrapozisyon}

``Yağmur yağarsa yer ıslanır'' doğruysa, ``yer ıslak değilse yağmur
yağmamıştır'' da doğrudur. Bu, \emph{kontrapozisyon} ilkesidir:
\[
  (P \to Q) \to (\neg Q \to \neg P).
\]

\leandosya[Kontrapozisyon]{Bolum03}{kontrapozisyon}

\noindent \leantac{intro hpq hnq hp} adımında \emph{üç} varsayım
içeri alınıyor. Üçüncüsü, sonuçtaki \lc{¬P}'nin içinde saklı olan
\lc{P}'dir. Bu adımdan sonra durum şöyledir:

\begin{hedefdurumu}[\texttt{intro hpq hnq hp} sonrası]
P Q : Prop
hpq : P → Q
hnq : ¬Q
hp : P
⊢ False
\end{hedefdurumu}

\noindent Terim modunda kontrapozisyon, yine bir fonksiyon bileşkesidir:
\lc{hpq}'yu uygulayıp sonucu \lc{hnq}'ya veririz.

\leandosya[Kontrapozisyon, terim modu]{Bolum03}{kontrapozisyon_terim}

\begin{figure}[htbp]
\centering
\begin{tikzpicture}[
    kanit/.style={draw=leanframe, fill=leanbg, thick, circle,
                  minimum size=0.95cm, font=\bfseries\color{leanframe}},
    makine/.style={draw=leanhyp, fill=leanhyp!12, thick, rounded corners=3pt,
                   minimum width=1.5cm, minimum height=0.9cm,
                   font=\leanfont\bfseries\color{leanhyp}},
    bos/.style={draw=hataframe, fill=hatabg, thick, circle,
                minimum size=0.95cm, font=\bfseries\color{hataframe}},
    ok/.style={-{Stealth[length=3mm]}, very thick, leansym},
  ]
  \node[kanit] (p) {$P$};
  \node[makine, right=0.9cm of p] (f) {hpq};
  \node[kanit, right=0.9cm of f] (q) {$Q$};
  \node[makine, right=0.9cm of q] (g) {hnq};
  \node[bos, right=0.9cm of g] (b) {$\bot$};
  \draw[ok] (p) -- (f); \draw[ok] (f) -- (q);
  \draw[ok] (q) -- (g); \draw[ok] (g) -- (b);
  \draw[decorate, decoration={brace, amplitude=6pt, mirror}, thick, vurgu]
    ([yshift=-0.75cm]p.west) -- ([yshift=-0.75cm]b.east)
    node[midway, below=9pt, font=\leanfont\small, text=vurgu]
    {fun hp => hnq (hpq hp)  :  P → False  =  ¬P};
\end{tikzpicture}
\caption{Kontrapozisyon: $P$'den $\bot$'a giden yol, $\neg P$'nin kanıtıdır.}
\label{fig:kontrapozisyon}
\end{figure}

Kontrapozisyonun \emph{tersi}, yani $(\neg Q \to \neg P) \to (P \to Q)$,
klasik mantıkta doğrudur ama yapıcı olarak ispatlanamaz. Bunu
\cref{bol:yapici-klasik}'te ispatlayacağız.

\leandosya[Modus tollens]{Bolum03}{modus_tollens}

\noindent Burada daha önce ispatladığımız bir teoremi (\lc{kontrapozisyon})
bir fonksiyon gibi çağırdık. Teoremler, bir kez ispatlandıktan sonra
yeniden kullanılabilen fonksiyonlardır.

% ---------------------------------------------------------------------
\section{Denklik (\texorpdfstring{$\leftrightarrow$}{↔})}\label{sec:denklik}

``$P$ ancak ve ancak $Q$'' önermesi, $P \to Q$ ile $Q \to P$'nin
birlikte doğru olduğunu söyler. \Lean{}'de \lc{P ↔ Q}, bu iki implikasyonun
bir çiftidir; dolayısıyla kuralları \lc{∧}'nınkilere çok benzer.

\begin{center}
\begin{tabular}{c@{\hspace{3em}}c@{\hspace{2em}}c}
\begin{prooftree}
  \hypo{[P]^1}
  \ellipsis{}{Q}
  \hypo{[Q]^2}
  \ellipsis{}{P}
  \infer2[$\leftrightarrow$G$^{1,2}$]{P \leftrightarrow Q}
\end{prooftree}
&
\begin{prooftree}
  \hypo{P \leftrightarrow Q}
  \hypo{P}
  \infer2[$\leftrightarrow$E$_1$]{Q}
\end{prooftree}
&
\begin{prooftree}
  \hypo{P \leftrightarrow Q}
  \hypo{Q}
  \infer2[$\leftrightarrow$E$_2$]{P}
\end{prooftree}
\\[10pt]
\small\textcolor{ikincil}{\bfseries giriş:} \leantac{constructor}, \lc{⟨f, g⟩} &
\multicolumn{2}{c}{\small\textcolor{vurgu}{\bfseries eleme:} \lc{h.mp}, \lc{h.mpr}, \leantac{rw [h]}}
\end{tabular}
\end{center}

\leandosya[Denkliğe giriş]{Bolum03}{iff_giris}
\leandosya[Denkliğe giriş, terim modu]{Bolum03}{iff_giris_terim}
\leandosya[Denkliğin elemesi]{Bolum03}{iff_eleme}

\noindent \lc{mp} adı \emph{modus ponens}'ten gelir; \lc{mpr} ise ``ters
yönde modus ponens'' demektir.

\subsection{Yeniden yazma: \texorpdfstring{\leantac{rw}}{rw}}

Denk önermeler, \emph{her bağlamda} birbirinin yerine kullanılabilir.
\leantac{rw [h]} taktiği, \lc{h : P ↔ Q} denkliğini kullanarak hedefteki
her \lc{P}'yi \lc{Q} ile değiştirir:

\leandosya[Yeniden yazma]{Bolum03}{iff_rw}

\begin{hedefgecisi}
\begin{hedefdurumu}[\texttt{rw [h]} öncesi]
h : P ↔ Q
hqr : Q → R
⊢ P → R
\end{hedefdurumu}
\sonra
\begin{hedefdurumu}[\texttt{rw [h]} sonrası]
h : P ↔ Q
hqr : Q → R
⊢ Q → R
\end{hedefdurumu}
\end{hedefgecisi}

\noindent \leantac{rw} hem eşitliklerle (\lc{a = b}) hem denkliklerle
(\lc{P ↔ Q}) çalışır. Eşitlikle yeniden yazmayı \cref{bol:yuklemler}'da
ayrıntılı olarak göreceğiz. Denklikle yeniden yazma ise \Lean{}'in
\emph{önerme dışsallığı} (\lc{propext}) aksiyomuna dayanır: Denk iki önerme
\emph{eşittir}. Bu aksiyoma \cref{bol:aksiyomlar}'da döneceğiz.

\subsection{Üçlü değilleme}

İki değilleme bir taneye indirgenemez, ama üç değilleme bire indirgenir:

\leandosya[Üçlü değilleme]{Bolum03}{uclu_degilleme}

\noindent Bu ispat, daha önce ispatladığımız \lc{cift\_degilleme\_giris}
teoremini yeniden kullanıyor. Sonuç şaşırtıcıdır: Yapıcı mantıkta
değillemeler dizisi $\neg\neg\neg\neg\neg P$ gibi ne kadar uzarsa uzasın,
yalnızca iki farklı önermeye indirgenir: $\neg P$ ya da $\neg\neg P$
(ve elbette $P$'nin kendisi).

% ---------------------------------------------------------------------
\section{De Morgan yasaları: yapıcılığın sınırı}\label{sec:demorgan}

Augustus De Morgan'ın adını taşıyan yasalar, değillemenin ve ile veya
üzerine nasıl dağıldığını söyler. Klasik mantıkta dört yönün hepsi
doğrudur:
\[
  \neg(P \vee Q) \leftrightarrow \neg P \wedge \neg Q
  \qquad\qquad
  \neg(P \wedge Q) \leftrightarrow \neg P \vee \neg Q
\]

\leandosya[De Morgan: veya'nın değillemesi]{Bolum03}{demorgan_veya}
\leandosya[De Morgan: ve'nin değillemesi, kolay yön]{Bolum03}{demorgan_ve_kolay}

\noindent Dördüncü yönü denediğimizde ise takılırız:

\begin{leankod}[Takıldığımız yer (derlenmez)]
theorem demorgan_ve_zor (P Q : Prop) : ¬(P ∧ Q) → ¬P ∨ ¬Q := by
  intro h
  -- h : ¬(P ∧ Q),  hedef: ⊢ ¬P ∨ ¬Q
  -- Or.inl mi, Or.inr mi? Hangisini seçeceğimizi bilmiyoruz!
  sorry
\end{leankod}

\begin{figure}[htbp]
\centering
\begin{tikzpicture}[
    yon/.style={draw=#1, fill=#1!8, thick, rounded corners=4pt,
                minimum width=3.7cm, minimum height=0.85cm, align=center,
                font=\small},
    ok/.style={-{Stealth[length=2.6mm]}, very thick, #1},
  ]
  \node[yon=anarenk] (a1) at (0,0)   {$\neg(P \vee Q)$};
  \node[yon=anarenk] (a2) at (7,0)   {$\neg P \wedge \neg Q$};
  \node[yon=anarenk] (b1) at (0,-2.1){$\neg(P \wedge Q)$};
  \node[yon=anarenk] (b2) at (7,-2.1){$\neg P \vee \neg Q$};

  \draw[ok=bittiframe] ([yshift=4pt]a1.east) -- node[above, font=\baslikfont\scriptsize\bfseries, text=bittiframe]{yapıcı {\sembolfont ✓}} ([yshift=4pt]a2.west);
  \draw[ok=bittiframe] ([yshift=-4pt]a2.west) -- node[below, font=\baslikfont\scriptsize\bfseries, text=bittiframe]{yapıcı {\sembolfont ✓}} ([yshift=-4pt]a1.east);
  \draw[ok=bittiframe] ([yshift=-4pt]b2.west) -- node[below, font=\baslikfont\scriptsize\bfseries, text=bittiframe]{yapıcı {\sembolfont ✓}} ([yshift=-4pt]b1.east);
  \draw[ok=hataframe, dashed] ([yshift=4pt]b1.east) -- node[above, font=\baslikfont\scriptsize\bfseries, text=hataframe]{yalnızca klasik} ([yshift=4pt]b2.west);
\end{tikzpicture}
\caption{De Morgan yasalarının dört yönü. Kesikli ok, yalnızca klasik
mantıkla ispatlanabilen yöndür.}
\label{fig:demorgan}
\end{figure}

Sorun şudur: Hedef bir \emph{veya} önermesidir ve onu kanıtlamanın tek
yolu (\cref{sec:veya}), iki taraftan \emph{birinin} kanıtını vermektir.
Elimizdeki tek bilgi ise ``$P$ ve $Q$ birlikte doğru olamaz''dır. Bu
bilgi, \emph{hangisinin} yanlış olduğunu söylemez. Gerçek hayattan bir
örnek:

\begin{gunlukornek}
``Bu bilgisayar programı hem hızlı hem doğru olamaz'' dediğimizi
düşünelim. Bundan ``program ya yavaştır ya da yanlıştır'' sonucu çıkar
mı? Klasik mantıkçı ``elbette'' der. Yapıcı mantıkçı ise şunu sorar:
\emph{Hangisi?} Programın yavaş olduğunu mu kanıtlayacaksın, yoksa yanlış
olduğunu mu? İkisinden birini gösteremiyorsan, ``ya o ya bu'' demeye
hakkın yoktur.
\end{gunlukornek}

Bu ``hangisi?'' sorusu, yapıcı mantığın özüdür. Yine de elimiz boş
değildir: Çift değilleme altında De Morgan'ın zor yönü yapıcı olarak
ispatlanabilir (Alıştırma~\ref{alis:demorgan-cift}). Bir sonraki kesim,
bu olgunun çok daha genel bir örneğini gösteriyor.

% ---------------------------------------------------------------------
\section{Üçüncü hâlin olanaksızlığı yanlış olamaz}\label{sec:em-yanlis-olamaz}

Klasik mantığın en ünlü ilkesi, her önermenin ya doğru ya yanlış
olduğunu söyleyen \emph{üçüncü hâlin olanaksızlığıdır} (Latince
\emph{tertium non datur}): $P \vee \neg P$. Bu önermeyi de yapıcı olarak
ispatlayamayız: $P$'yi mi yoksa $\neg P$'yi mi kanıtlayacağımızı bilmeden
\lc{Or.inl} ile \lc{Or.inr} arasında seçim yapamayız. Ama şaşırtıcı bir
şey yapabiliriz: Bu ilkenin \emph{yanlış olamayacağını} ispatlayabiliriz.

\leandosya[\texorpdfstring{$\neg\neg(P \vee \neg P)$}{¬¬(P ∨ ¬P)}]{Bolum03}{em_cift_degil}

Bu ispat, mantıkta bir tür ``hile'' içerir ve adım adım izlenmeye değer:

\begin{enumerate}[adimlar]
  \item \lc{h : ¬(P ∨ ¬P)} varsayıyoruz ve bir çelişki arıyoruz.
  \item \leantac{apply h}: Çelişki elde etmenin bir yolu, \lc{h}'yi
        çürütmek, yani \lc{P ∨ ¬P}'nin bir kanıtını vermek.
  \item \leantac{apply Or.inr}: \emph{Cesur bir seçim yapıyoruz}:
        $\neg P$ tarafını kanıtlayacağız.
  \item \leantac{intro hp}: $\neg P$'yi kanıtlamak için $P$'yi varsayıp
        çelişki arıyoruz.
  \item \leantac{apply h}, \leantac{exact Or.inl hp}: Ama artık elimizde
        $P$ var! Seçimimizi değiştirip \lc{P ∨ ¬P}'nin \emph{sol}
        tarafını kanıtlıyor ve \lc{h} ile çelişki elde ediyoruz.
\end{enumerate}

\begin{gunlukornek}[Bir benzetme: şeytanın teklifi]
Bu ispatın programlama karşılığı, mantıkçılar arasında ``şeytanın
teklifi'' olarak anlatılan bir hikâyeye benzer. Şeytan der ki: ``Ya sana
bir milyar lira vereceğim, ya da bir milyar lira ödersen dilediğin
her şeyi yapacağım. Hangisini yapacağıma ben karar veririm.'' Ve ikinci
seçeneği söyler. Yıllar sonra sen bir milyar lirayı biriktirip getirdiğinde,
şeytan ``Fikrimi değiştirdim: birinci seçenek. İşte paran'' der ve az önce
verdiğin parayı sana geri verir. İspattaki \lc{Or.inr}'den \lc{Or.inl}'ye
geçiş tam olarak budur: Önce $\neg P$ diyor, $P$'nin kanıtı gelince
o kanıtı kullanarak ``aslında $P$'' diyoruz.
\end{gunlukornek}

\leandosya[\texorpdfstring{$\neg\neg(P \vee \neg P)$}{¬¬(P ∨ ¬P)}, terim modu]{Bolum03}{em_cift_degil_terim}

Elimizde artık iki gözlem var: \lc{P ∨ ¬P}'yi yapıcı olarak
ispatlayamıyoruz, ama \lc{¬¬(P ∨ ¬P)}'yi ispatlayabiliyoruz. Arada kalan
tek adım $\neg\neg X \to X$, yani \emph{çift değilleme eleme}dir. Bu adım,
yapıcı mantık ile klasik mantığı ayıran sınırın kendisidir.
\Cref{bol:yapici-klasik} tamamen bu sınıra ayrılmıştır.

% ---------------------------------------------------------------------
\section{Curry--Howard tablosuna eklemeler}\label{sec:curry-howard-2}

\begin{center}
\renewcommand{\arraystretch}{1.5}
\rowcolors{2}{tablosatir}{white}
\begin{tabular}{>{\raggedright\arraybackslash}p{2.7cm}
                >{\raggedright\arraybackslash}p{4.0cm}
                >{\raggedright\arraybackslash}p{5.0cm}}
\rowcolor{tabloust}
\tablobaslik{Mantık} & \tablobaslik{Programlama} & \tablobaslik{Lean'deki kanıt} \\
$\top$              & tek elemanlı tip (\emph{unit}) & \lc{True.intro}, \leantac{trivial} \\
$\bot$              & boş tip                         & --- (kanıt yoktur) \\
$\bot$ eleme        & boş tipten fonksiyon            & \lc{False.elim h}, \leantac{cases h} \\
$\neg P$            & $P$'den boş tipe fonksiyon      & \lc{fun hp => ...} (\lc{False} üretir) \\
$P \leftrightarrow Q$ & iki yönlü dönüşüm çifti       & \lc{⟨f, g⟩} \\
\end{tabular}
\end{center}

% ---------------------------------------------------------------------
\section{Taktik özeti}

\begin{center}
\renewcommand{\arraystretch}{1.45}
\rowcolors{2}{tablosatir}{white}
\begin{tabularx}{\linewidth}{>{\raggedright\arraybackslash}p{3.1cm}
                             >{\raggedright\arraybackslash}p{1.8cm}
                             >{\raggedright\arraybackslash}X}
\rowcolor{tabloust}
\tablobaslik{Taktik / terim} & \tablobaslik{Kural} & \tablobaslik{Ne yapar?} \\
\leantac{trivial}               & $\top$G & \lc{⊢ True} hedefini kapatır. \\
\leantac{exfalso}               & $\bot$E & Hedefi \lc{False}'a çevirir. \\
\lc{False.elim h}               & $\bot$E & \lc{h : False}'tan her şeyi üretir. \\
\lc{absurd hp hnp}              & $\neg$E, $\bot$E & \lc{P} ve \lc{¬P}'den her şeyi üretir. \\
\leantac{contradiction}         & ---     & Bağlamdaki açık çelişkiyi bulur. \\
\leantac{intro hp} (hedef \lc{¬P}) & $\neg$G & \lc{hp : P} varsayar, hedef \lc{False} olur. \\
\leantac{constructor} (hedef \lc{↔}) & $\leftrightarrow$G & İki yönü ayrı hedefler olarak açar. \\
\lc{h.mp}, \lc{h.mpr}           & $\leftrightarrow$E & Denkliği bir yönde kullanır. \\
\leantac{rw [h]}                & ---     & \lc{h : P ↔ Q} ile hedefteki \lc{P}'leri \lc{Q} yapar. \\
\end{tabularx}
\end{center}

% ---------------------------------------------------------------------
\begin{arastirmanotu}
\textbf{Değillemenin tanımı.} $\neg P := P \to \bot$ tanımı, Heyting'in
1930'da sezgici mantığı biçimselleştirirken yaptığı seçime dayanır
\citep{heyting1930}. Bu tanım, değillemeyi bağımsız bir kavram olmaktan
çıkarıp implikasyona indirger ve Curry--Howard eşlemesinin
değillemeye de uzanmasını sağlar \citep{sorensen2006}.

\textbf{İki değilleme sınırı.} Yapıcı önermeler mantığında
$\neg\neg\neg P \leftrightarrow \neg P$ olduğundan, tek bir önerme
değişkeninden yalnızca değilleme ile sonlu sayıda farklı önerme elde
edilir. Tek değişkenli önermelerin tamamı ise sonsuz bir yapı oluşturur;
bu yapı \emph{Rieger--Nishimura kafesi} olarak bilinir.

\textbf{Paraconsistent mantıklar.} Ex falso ilkesi bütün mantıklarda
kabul edilmez. \emph{Paraconsistent} (çelişkiye dayanıklı) mantıklar,
çelişkili bilgiyle çalışan sistemlerde (ör. tutarsız veri tabanları)
her şeyin türetilebilmesini engellemek için bu ilkeyi bırakır. Bu
mantıklar kitabın kapsamı dışındadır.
\end{arastirmanotu}

\begin{bolumozeti}
\begin{itemize}
  \item \lc{True} tek elemanlı, \lc{False} boş tiptir. \lc{False}'tan her
        şey çıkar (\emph{ex falso}).
  \item \lc{¬P}, \lc{P → False}'un kısaltmasıdır. Değillemeyi kanıtlamak
        için \leantac{intro}, kullanmak için fonksiyon uygulaması yeter.
  \item Çelişki üç yoldan kullanılır: \leantac{exfalso}, \lc{absurd},
        \leantac{contradiction}.
  \item \lc{P ↔ Q} iki implikasyonun çiftidir; \lc{h.mp}, \lc{h.mpr} ve
        \leantac{rw} ile kullanılır.
  \item Yapıcı mantıkta $P \to \neg\neg P$ doğrudur, tersi değil. De
        Morgan'ın bir yönü, $P \vee \neg P$ ve $\neg\neg P \to P$
        yapıcı olarak ispatlanamaz; ama $\neg\neg(P \vee \neg P)$
        ispatlanabilir.
\end{itemize}
\end{bolumozeti}

\begin{alistirmalar}
\alistirma{1} \lc{¬False} önermesini ispatlayın. Bu önermenin
\lc{False → False}'tan farkı var mı? (\lc{a3\_1})

\alistirma{1} \lc{P → True} önermesini ispatlayın. (\lc{a3\_2})

\alistirma{1} \lc{P → ¬P → Q} önermesini önce \leantac{exfalso}, sonra
\lc{absurd} ile ispatlayın. (\lc{a3\_3})

\alistirma{1} Denkliğin yansıma özelliğini ispatlayın: \lc{P ↔ P}.
(\lc{a3\_4})

\alistirma{2} Denkliğin simetrisini (\lc{a3\_5}) ve geçişliliğini
(\lc{a3\_6}) ispatlayın.

\alistirma{2} \emph{Disjunktif kıyas}: \lc{P ∨ Q → ¬P → Q}.
``Ya çay ya kahve içeceğim; çay içmeyeceğim; o hâlde kahve içeceğim.''
(\lc{a3\_7})

\alistirma{2} Hiçbir önerme kendi değillemesine denk olamaz:
\lc{¬(P ↔ ¬P)}. İpucu: Önce \leantac{have} ile \lc{¬P}'yi elde edin.
Bu alıştırmanın \emph{yalancı paradoksu} (``Bu cümle yanlıştır'') ile
ilişkisini tartışın. (\lc{a3\_8})

\alistirma{2} \lc{¬(P → Q) → ¬Q} önermesini ispatlayın. (\lc{a3\_9})

\alistirma{2} Çift değillemenin implikasyona dağıldığını ispatlayın:
\lc{¬¬(P → Q) → ¬¬P → ¬¬Q}. (\lc{a3\_10})

\alistirma{3}\label{alis:demorgan-cift} De Morgan'ın zor yönünün çift
değilleme altında yapıcı olduğunu ispatlayın:
\lc{¬(P ∧ Q) → ¬¬(¬P ∨ ¬Q)}. İpucu: \lc{em\_yanlis\_olamaz}
ispatındaki ``fikir değiştirme'' hilesini kullanın. (\lc{a3\_11})

\alistirma{3}\label{alis:peirce-cift} Bölüm~2'deki Peirce yasasının
çift değilleme altında yapıcı olduğunu ispatlayın:
\lc{¬¬(((P → Q) → P) → P)}. (\lc{a3\_12})

\alistirma{3} (Kâğıt-kalem) \lc{em\_yanlis\_olamaz} ispatının doğal
çıkarım ağacını çizin. \lc{h} varsayımı ağaçta kaç kez kullanılıyor?
Bu ispatı, varsayımı yalnızca bir kez kullanan bir ispatla
değiştirmek mümkün müdür? Tartışın.
\end{alistirmalar}
